\documentclass[journal]{IEEEtran}
\usepackage{amsmath,amssymb,bm}
\usepackage{amsthm}
\usepackage{graphicx}
\usepackage[caption=false,font=footnotesize]{subfig}
\usepackage{booktabs}
\usepackage{cite}
\usepackage[hidelinks]{hyperref}
\usepackage{url}
\theoremstyle{definition}
\newtheorem{assumption}{Assumption}
\newtheorem{definition}{Definition}
\newtheorem{problem}{Problem}
\theoremstyle{plain}
\newtheorem{theorem}{Theorem}
\newtheorem{lemma}{Lemma}
\newtheorem{corollary}{Corollary}
\theoremstyle{remark}
\newtheorem{remark}{Remark}

% IEEE TAES Correspondence version.
% Target: <= 8 pages (two-column IEEEtran), all figures/tables embedded.
% Long derivations are confined to the Appendix, which can be moved to the
% supplementary material if a shorter limit applies.

\title{Prescribed-Time Joint Tracking of a Free-Floating Space Manipulator via Smooth Periodic Delayed Feedback}

\author{Author~Name and Author~Name\\
\thanks{Manuscript received XX Month 2025; revised XX Month 2025; accepted XX Month 2025. Date of publication XX Month 2025; date of current version XX Month 2025. This work was supported in part by XXX under Grant XXX. \emph{(Corresponding author: XXX.)}
The authors are with the Department of XXX, XXX University, XXX 000000, China (e-mail: author@xxx.edu.cn).}}


\begin{document}
\maketitle

\begin{abstract}
Free-floating space manipulators are strongly coupled, torque limited, and their settling time is normally fixed by the controller gains rather than by the task. This correspondence shows that such a manipulator can be converted into a controllable error channel for smooth periodic delayed feedback (PDF), so that a designer-specified settling time is met without nonsmooth terminal terms. Momentum conservation and the Schur complement of the full spatial inertia matrix give an exact joint-space model, a prescribed-time observer with a designable observation time provides the disturbance interface, and a Gramian-based periodic gain annihilates the one-period state-transition operator of the error channel. In the ideal case the tracking error is exactly zero after $T_o+2h$; with observation error, model mismatch, actuator limits, and sampling, an explicit practical bound is obtained. On the three-dimensional model of a seven-degree-of-freedom space manipulator with tabulated parameters and $\pm20~\mathrm{N\,m}$ actuators, the delayed feedback reduces the post-settling error from $1.4\times10^{-3}$ to $4.3\times10^{-5}\,\mathrm{rad}$, and the observer reduces the disturbed error from $2.6\times10^{-1}$ to $1.5\times10^{-3}\,\mathrm{rad}$ without saturating the actuators. A design rule relating the inertia spread to the Gramian conditioning is reported.
\end{abstract}

\begin{IEEEkeywords}
Delay feedback, disturbance observers, free-floating manipulators, prescribed-time control, trajectory tracking.
\end{IEEEkeywords}

\section{Introduction}
\IEEEPARstart{F}{ree-floating} space manipulators are the enabling technology for on-orbit servicing, debris removal, and in-space assembly~\cite{liu2021review,yoshida2003etsvii}. Before capture, the base is not rigidly supported: joint motion produces a reaction motion of the base through momentum conservation, so the end-effector is coupled to the entire spacecraft--manipulator system~\cite{dubowsky1993,papadopoulos1991nature}. This coupling, together with parameter uncertainty, disturbances, and actuator limits, makes trajectory tracking substantially harder than for a fixed-base arm, and the difficulty is aggravated when the settling time is a task requirement rather than a by-product of the gains.

\emph{Related work.} Three lines are relevant. First, the coupled kinematics is described by the generalized Jacobian~\cite{umetani1989}, whose configuration-dependent singularity loci are now well understood~\cite{papadopoulos1993,nanos2015avoiding}, and the base motion can be eliminated to obtain an equivalent joint-space model~\cite{parlaktuna2004}. Second, robust tracking has been addressed by disturbance observers~\cite{chu2008dob}, adaptive Jacobian schemes~\cite{wang2012adaptivejacobian}, time-delay estimation~\cite{wang2025tdelog}, and terminal sliding modes with finite-time observers~\cite{shao2021ntsm,dou2023fractional,dou2025fractional}. Third, the convergence time itself has become a design objective: finite-time~\cite{jia2020finite}, fixed-time~\cite{polyakov2012,hu2024constrainedfixed}, and predefined- or prescribed-time designs~\cite{xu2024se3predefined,zou2025practical,yan} specify it explicitly. Most of the last group relies on fractional powers, terminal nonsmooth terms, or gains that grow without bound near the terminal instant, and the disturbance observers they use offer no handle on the observation time~\cite{yan}. Saturation and performance envelopes add further conservatism when several constraints coexist~\cite{zhou2025saturated}.

\emph{Delayed feedback.} Smooth periodic delayed feedback (PDF) is an alternative that has no terminal singularity: for a controllable linear time-delay system, a periodic delayed gain annihilates the one-period state-transition operator, which yields fixed-time stabilisation with a continuous input~\cite{pyragas1992,zhouFixedtimeStabilizationLinear2022}. Extending this mechanism to a free-floating manipulator requires three ingredients that are not available in~\cite{zhouFixedtimeStabilizationLinear2022}: an exact auxiliary input channel that absorbs the nonlinear and momentum-coupled terms; a disturbance interface whose observation time is designable, for which a prescribed-time observer is the natural candidate~\cite{jiangPrescribedtimeDisturbanceObserver2024}; and a quantification of how the residual observation error, saturation, and sampling perturb the one-period map.

\emph{Contributions.} This correspondence
\begin{enumerate}
\item reduces the free-floating manipulator to an exact controllable error channel using the Schur complement of the spatial inertia matrix and the momentum-consistent bias forces, and designs a smooth PDF law on that channel;
\item states the one-period nullification theorem for the reduced nonlinear plant, together with a practical bound in which the observation residual, the actuator limit, and the discretisation appear explicitly;
\item reports a design rule for the Gramian gain: because the equivalent inertia of a long-reach arm spans three decades, normalising the gains by the joint inertia spreads the modal time scales and makes the Gramian numerically singular, whereas a uniform gain keeps the Gramian block diagonal and well conditioned;
\item evaluates the design on a complete three-dimensional model of a seven-degree-of-freedom manipulator with published parameters and $\pm20~\mathrm{N\,m}$ actuators, reporting settling time, peak torque, saturation duty cycle, and base reaction motion.
\end{enumerate}

\emph{Organization.} Section~\ref{sec:model} states the reduced model, Section~\ref{sec:control} the observer and the PDF law, Section~\ref{sec:stability} the stability and practical bound, Section~\ref{sec:simulation} the numerical example, and Section~\ref{sec:conclusion} the conclusions. Long derivations are collected in the Appendix.

\section{Reduced Dynamics}\label{sec:model}
Consider a free-floating manipulator with a rigid base and $n$ revolute joints, with joint vector $q\in\mathbb R^n$, joint torque $\tau\in\mathbb R^n$, base velocity $\dot x_b=[v_b^\top,\omega_b^\top]^\top\in\mathbb R^6$, and end-effector velocity $\dot x_e\in\mathbb R^6$.

\subsection{Momentum reduction}
With no external force or torque, the kinetic energy is
\begin{equation}
 T=\tfrac12\begin{bmatrix}\dot x_b^\top&\dot q^\top\end{bmatrix}
 \begin{bmatrix}H_b&H_{bm}\\H_{bm}^\top&H_m\end{bmatrix}
 \begin{bmatrix}\dot x_b\\\dot q\end{bmatrix},
 \label{eq:kinetic}
\end{equation}
and the Lagrange equations read
\begin{equation}
 \begin{bmatrix}H_b&H_{bm}\\H_{bm}^\top&H_m\end{bmatrix}
 \begin{bmatrix}\ddot x_b\\\ddot q\end{bmatrix}
 +\begin{bmatrix}C_b&C_{bm}\\C_{bm}^\top&C_m\end{bmatrix}
 \begin{bmatrix}\dot x_b\\\dot q\end{bmatrix}
 =\begin{bmatrix}0\\\tau\end{bmatrix}.
 \label{eq:full}
\end{equation}

\begin{assumption}\label{as:momentum}
The initial linear and angular momenta are zero, i.e.,
$H_b(q)\dot x_b+H_{bm}(q)\dot q=0$, and $H_b(q)$ is uniformly positive
definite on the workspace.
\end{assumption}

Under Assumption~\ref{as:momentum},
\begin{equation}
 \dot x_b=-H_b^{-1}(q)H_{bm}(q)\dot q\triangleq J_{bm}(q)\dot q,
 \label{eq:jbm}
\end{equation}
\begin{equation}
 \dot x_e=\bigl[J_m(q)+J_b(q)J_{bm}(q)\bigr]\dot q\triangleq J_g(q)\dot q,
 \label{eq:gjm}
\end{equation}
where $J_g$ is the generalized Jacobian. Differentiating the momentum
constraint and eliminating $\ddot x_b$ gives the exact equivalent model
\begin{equation}
 M_e(q)\ddot q+C_e(q,\dot q)\dot q=\tau+d(t),
 \label{eq:reduced}
\end{equation}
with
\begin{equation}
\begin{aligned}
 M_e&=H_m-H_{bm}^\top H_b^{-1}H_{bm},\\
 C_e\dot q&=C_m-H_{bm}^\top H_b^{-1}C_b,
\end{aligned}
 \label{eq:schur}
\end{equation}
where $C_b$ and $C_m$ are the bias forces of \eqref{eq:full} evaluated at the
momentum-consistent base velocity \eqref{eq:gjm}. The signal $d(t)$ collects
parameter errors, unmodeled dynamics, and external disturbances.

\begin{assumption}\label{as:inertia}
There is $\underline m>0$ such that $M_e(q)\succeq\underline m I_n$ on the
workspace, and the nominal $M_e$, $C_e$ are available to the controller.
\end{assumption}

\begin{remark}\label{rem:implementation}
Equations \eqref{eq:schur} are implemented exactly: $H_{\mathrm{full}}$ is
assembled with the spatial composite rigid-body algorithm, the bias forces with
a recursive Newton--Euler pass at $\dot x_b=J_{bm}\dot q$, and the Schur
complement and $C_e\dot q$ are extracted afterwards. For the seven-degree-of-freedom
manipulator of Section~\ref{sec:simulation} the implementation returns
$\lambda_{\min}(H_{\mathrm{full}})=0.1375$, a momentum residual
$\|H_b\dot x_b+H_{bm}\dot q\|=4.2\times10^{-15}$, a skew-symmetry residual
$|\dot q^\top(\dot M_e\dot q-2C_e\dot q)|=2.0\times10^{-10}$, and a relative
kinetic-energy drift of $3.3\times10^{-6}$ over one second with $\tau=0$.
\end{remark}

\begin{remark}\label{rem:spread}
For the same manipulator, $\mathrm{diag}(M_e)$ varies between $0.14$ and
$380~\mathrm{kg\,m^2}$, i.e., over a factor of $2.7\times10^3$. This spread,
not the control law, fixes the numerical conditioning of the design in
Section~\ref{sec:control} and the admissible initial error for a given torque
limit.
\end{remark}

\section{Controller Design}\label{sec:control}

\subsection{Tracking objective and timing}
Let $q_d\in C^2([0,\infty);\mathbb R^n)$ be a bounded reference with bounded
derivatives and define
\begin{equation}
 e=q-q_d,\qquad \dot e=\dot q-\dot q_d,\qquad z=[e^\top,\dot e^\top]^\top .
 \label{eq:errors}
\end{equation}
The total settling time is prescribed by two design parameters, the
observation time $T_o$ and the delay $h$:
\begin{equation}
 T=T_o+2h .
 \label{eq:prescribed}
\end{equation}
The objective is $e(t)=\dot e(t)=0$ for $t\ge T$ under exact compensation,
and an explicit ultimate bound in terms of the residual acceleration otherwise.

\subsection{Prescribed-time disturbance observer}\label{sec:ptdo}
Let $\chi=\dot q$ and write the mechanical subsystem as
\begin{equation}
 \dot\chi=u_o+\Delta_a,\qquad
 u_o=M_e^{-1}\bigl[\tau-C_e\dot q\bigr],\qquad
 \Delta_a=M_e^{-1}d .
 \label{eq:channel}
\end{equation}
The following observer is used:
\begin{equation}
 \dot z_1=z_2+u_o+\ell_1(t)\phi_1(\varepsilon_1/\sigma),\qquad
 \dot z_2=\ell_2(t)\phi_2(\varepsilon_1/\sigma),
 \label{eq:ptdo}
\end{equation}
with $\varepsilon_1=\chi-z_1$,
\begin{equation}
\begin{aligned}
 \phi_1(x)&=\operatorname{sig}^{1-\eta/2}(x)+\operatorname{sig}^{1+\eta/2}(x),\\
 \phi_2(x)&=\operatorname{sig}^{1-\eta}(x)+\operatorname{sig}^{1+\eta}(x),
\end{aligned}
 \label{eq:phi}
\end{equation}
$\operatorname{sig}^{\alpha}(x)=|x|^{\alpha}\operatorname{sgn}(x)$,
$0<\eta<1$, and the prescribed-time gains
\begin{equation}
 \ell_1(t)=\frac{m_1}{T_o-t},\qquad
 \ell_2(t)=\frac{m_2}{(T_o-t)^2},\qquad t<T_o .
 \label{eq:gains}
\end{equation}
The powers in \eqref{eq:phi} are dilation-consistent,
$1-\eta=2(1-\eta/2)-1$ and $1+\eta=2(1+\eta/2)-1$, so each branch of the error
system is homogeneous of negative degree with weights $(1,1\mp\eta/2)$; the
gains \eqref{eq:gains} serve as the time rescaling $s=-\ln(1-t/T_o)$.

\begin{assumption}[Observer interface]\label{as:interface}
There exist $T_o>0$ and an estimate $\hat\Delta_a$ such that, under the
continuous-time observer \eqref{eq:ptdo}, $\hat\Delta_a(t)=\Delta_a(t)$ for all
$t\ge T_o$, with a bounded error before $T_o$. The disturbance estimate fed to
the controller is $\hat d=M_e\hat\Delta_a$.
\end{assumption}

\begin{remark}\label{rem:frozen}
In the implementation the gains \eqref{eq:gains} are frozen at their value at
$T_o-\delta$. The frozen observer is a high-gain second-order differentiator
with residual $O(\|\dot\Delta_a\|\delta^2/m_2)$, which is part of the residual
$\Delta_r$ of Section~\ref{sec:stability}; it affects the size of the practical
bound, not the stability statement.
\end{remark}

\subsection{Dynamic compensation}
Using
\begin{equation}
 \tau=M_e(q)v+C_e(q,\dot q)\dot q-\hat d,
 \qquad v=\ddot q_d+\nu,
 \label{eq:computed_torque}
\end{equation}
and invoking Assumption~\ref{as:interface} for $t\ge T_o$, the error dynamics
of \eqref{eq:reduced} become the double integrator
\begin{equation}
 \dot z=Az+B\nu,\qquad
 A=\begin{bmatrix}0&I\\0&0\end{bmatrix},\quad
 B=\begin{bmatrix}0\\I\end{bmatrix},
 \label{eq:error_channel}
\end{equation}
which is controllable since $\mathrm{rank}[B,AB]=2n$.

\subsection{Smooth periodic delayed feedback}\label{sec:pdf}
Choose $K_0$ so that $A_c=A+BK_0$ is Hurwitz, let $h>0$, and let
$R_h:[0,2h]\to\mathbb R$ be smooth, nonnegative, zero on $[0,h]$ and positive
on a subset of positive measure in $(h,2h)$. Define the Gramian
\begin{equation}
 W_c=\int_h^{2h}e^{-A_cs}BR_h(s)B^\top e^{-A_c^\top s}\,ds
 \label{eq:wc}
\end{equation}
and, for the local time $\vartheta=t-T_o$, the $2h$-periodic gain
\begin{equation}
 K_c(\vartheta)=
 \begin{cases}
 0, & 0\le\vartheta<h,\\
 R_h(\vartheta)B^\top e^{-A_c^\top\vartheta}W_c^{-1}
 e^{A_c(h-\vartheta)}, & h\le\vartheta<2h .
 \end{cases}
 \label{eq:kc}
\end{equation}
The control input is
\begin{equation}
 \nu(t)=K_0z(t)-K_c(t-T_o)z(t-h),\qquad t\ge T_o,
 \label{eq:nu}
\end{equation}
so that the complete law is
\begin{equation}
 \tau=M_e\bigl[\ddot q_d+K_0z-K_c(t-T_o)z(t-h)\bigr]+C_e\dot q-\hat d .
 \label{eq:law}
\end{equation}
During the first half-period $K_c=0$ and the delayed state is only recorded;
the delayed term acts in the second half-period, where its product with the
phase-one trajectory cancels the period map.

\begin{remark}[Gain selection]\label{rem:gains}
A uniform baseline gain $K_0=[-K_p,\,-K_d]$ with $K_p=k_pI_n$,
$K_d=k_dI_n$ is used. Then $A_c$ is block diagonal with identical $2\times2$
blocks and $W_c$ inherits the same structure, with a condition number of
$7.4\times10^3$ for the example below. Normalising $K_p$ by the joint inertia
equalises the torque per radian of error but spreads the modal time scales over
three decades; in that case $\mathrm{cond}(W_c)$ exceeds $10^8$ and the
one-period nullification is destroyed by the inversion in \eqref{eq:kc}. The
uniform choice is therefore both numerically sound and, for this actuator
class, physically necessary.
\end{remark}

\section{Stability and Practical Bound}\label{sec:stability}

\begin{theorem}[One-period nullification]\label{thm:null}
Suppose Assumptions~\ref{as:momentum}--\ref{as:interface} hold, the model is
exact, the actuator is not saturated, and $W_c\succ0$. Then the error channel
\eqref{eq:error_channel} under \eqref{eq:nu} satisfies
$z(T_o+2h)=0$ for every finite $z(T_o)$, and $z(t)=0$ for all
$t\ge T_o+2h$.
\end{theorem}

\begin{proof}
For $0\le\vartheta<h$ the delayed term vanishes, hence
$z(T_o+\vartheta)=e^{A_c\vartheta}z(T_o)$ and, in particular,
$z(T_o+h)=e^{A_ch}z(T_o)$. For $h\le\vartheta\le2h$, variation of constants
gives
\begin{equation}
\begin{aligned}
 z(T_o+2h)=\Bigl[&e^{2A_ch}\\
 &-\int_h^{2h}e^{A_c(2h-s)}BK_c(s)e^{A_c(s-h)}ds\Bigr]z(T_o).
\end{aligned}
 \label{eq:voc}
\end{equation}
Substituting \eqref{eq:kc}, and using $e^{A_c(h-s)}e^{A_c(s-h)}=I$, the
integral equals
\begin{multline}
 \int_h^{2h}e^{A_c(2h-s)}BR_h(s)B^\top e^{-A_c^\top s}
 W_c^{-1}e^{A_c(h-s)}e^{A_c(s-h)}ds\\
 =e^{2A_ch}\Bigl[\int_h^{2h}e^{-A_cs}BR_h(s)B^\top e^{-A_c^\top s}ds\Bigr]W_c^{-1}
 =e^{2A_ch},
\end{multline}
so $z(T_o+2h)=0$. Since the delayed history is zero afterwards and the
dynamics are homogeneous, uniqueness gives $z(t)=0$ for all later time.
\end{proof}

\begin{corollary}\label{cor:exact}
Under the assumptions of Theorem~\ref{thm:null}, $e(t)=\dot e(t)=0$ for all
$t\ge T_o+2h=T$, and the settling time is independent of $z(0)$.
\end{corollary}

\begin{theorem}[Practical bound]\label{thm:practical}
Let the implemented model, the observation residual, and the saturation
generate a residual acceleration $\Delta_r$ in \eqref{eq:error_channel}, and
suppose $\|\Delta_r(t)\|\le\bar\Delta_r$ for $t\ge T_o$. Let
$\Psi(\vartheta,s)$ be the causal transition operator of the homogeneous
delayed system and
\begin{equation}
 \Gamma_h=\sup_{0\le\vartheta\le2h}\int_0^\vartheta\|\Psi(\vartheta,s)B\|\,ds .
 \label{eq:gamma}
\end{equation}
Then, for $t_k=T_o+2kh$,
\begin{equation}
\begin{aligned}
 \|z(t_{k+1})\|&\le\Gamma_h\bar\Delta_r,\qquad k\ge0,\\
 \|z(t)\|&\le\Gamma_h\bar\Delta_r,\qquad t\ge T_o+2h .
\end{aligned}
 \label{eq:bound}
\end{equation}
\end{theorem}

\begin{proof}
Variation of constants over one period gives
$z(t_{k+1})=\Delta_Kz(t_k)+\int_{t_k}^{t_{k+1}}\Psi(t_{k+1},s)B\Delta_r(s)ds$
with $\Delta_K=0$ by Theorem~\ref{thm:null}; taking norms and using
\eqref{eq:gamma} yields the boundary estimate, and the within-period estimate
follows because the homogeneous period-boundary state is zero and the
within-period transition operators are bounded by the same constant.
\end{proof}

If the actuator limits are active or the model is perturbed during the design
window, the one-period map is only approximately annihilated. Writing
$\|\Delta_K\|\le\varepsilon_\Delta<1$ gives the geometric bound
\begin{equation}
 \limsup_{k\to\infty}\|z(t_k)\|\le\frac{\Gamma_h\bar\Delta_r}
 {1-\varepsilon_\Delta},
 \label{eq:eps_bound}
\end{equation}
which is the form that applies to the saturated case of
Section~\ref{sec:simulation}. If $\|M_e^{-1}\|\le\underline m^{-1}$ and the
observation, model, and actuation residuals are bounded by $\bar d_o$,
$\bar\Delta_m$, $\bar\Delta_\tau$, then
\begin{equation}
 \bar\Delta_r\le\underline m^{-1}\bigl(\bar d_o+\bar\Delta_m+\bar\Delta_\tau\bigr).
 \label{eq:delta_r}
\end{equation}
The bound \eqref{eq:bound} is conservative: it is a worst-case estimate over
the period, whereas the measured residuals in Section~\ref{sec:simulation} are
two orders of magnitude smaller.

\section{Numerical Example}\label{sec:simulation}

\subsection{Plant and settings}
The manipulator has a rigid base and seven revolute joints. Its kinematic and
dynamic parameters, its initial configuration, and its $\pm20~\mathrm{N\,m}$
actuator limits are taken from a published study of the same system
class~\cite{dou2025fractional} and are listed in Table~\ref{tab:params}; no
parameter is tuned here. The chain is written in modified
Denavit--Hartenberg form
\begin{equation}
 A_i=\mathrm{R}_x(\alpha_i)\mathrm{T}_x(L_i)\mathrm{R}_z(\theta_i)\mathrm{T}_z(d_i),
 \qquad \theta_i=q_i+\theta_{o,i},
 \label{eq:mdh}
\end{equation}
so that joint $i$ rotates about the $z$ axis of link frame $i$ and the link
vector $L_i$ lies along its $x$ axis, the link centre of mass being at
$L_i/2$. This convention reproduces the tabulated data: the two long links
($L_4=L_5=1.2~\mathrm{m}$) have their centres of mass at $0.6~\mathrm{m}$ and
are slender along $x$, consistent with the tabulated
$I_{xx}\ll I_{yy}\simeq I_{zz}$. The base frame is attached to the spacecraft
centre of mass and the first joint is at $(0,0,0.72)~\mathrm{m}$ in that frame.
The base has mass $1000~\mathrm{kg}$ and inertia
$\mathrm{diag}(1200,1200,1200)~\mathrm{kg\,m^2}$; the initial joint angles are
$[0,-90,45,-45,90,90,0]^\circ$.

\begin{table}[!t]
\centering
\caption{Parameters of the 7-DOF manipulator, from~\cite{dou2025fractional}. Masses in kg, lengths in m, inertias in $\mathrm{kg\,m^2}$ about the link centre of mass.}
\label{tab:params}
\footnotesize
\setlength{\tabcolsep}{3pt}
\begin{tabular}{lccccccc}
\toprule
Link $i$ & 1 & 2 & 3 & 4 & 5 & 6 & 7\\
\midrule
$\alpha_i$ [deg] & 90 & 90 & $-90$ & 0 & 0 & 90 & $-90$\\
$L_i$ [m]        & 0 & 0 & 0 & 1.2 & 1.2 & 0 & 0\\
$d_i$ [m]        & 0.18 & 0.22 & 0.20 & 0.24 & 0.22 & 0.20 & 0.30\\
$\theta_{o,i}$ [deg] & 0 & 0 & $-90$ & 0 & 90 & 0 & 0\\
$m_i$ [kg]       & 35.01 & 30.00 & 22.69 & 21.38 & 16.75 & 26.17 & 18.10\\
$I_{xx,i}$   & 1.218 & 2.100 & 0.102 & 0.4327 & 0.3878 & 0.5727 & 0.170\\
$I_{yy,i}$   & 0.5132 & 1.378 & 3.378 & 2.266 & 0.3963 & 0.5987 & 0.240\\
$I_{zz,i}$   & 1.331 & 2.359 & 3.359 & 1.911 & 0.07271 & 0.1288 & 0.140\\
\bottomrule
\end{tabular}
\end{table}

The sampling step is $T_s=10^{-3}~\mathrm{s}$ and the horizon is $10~\mathrm{s}$.
The reference is generated by staggered fifth-order polynomials, joint $i$
completing its displacement of $[15,22,-18,22,-22,28,18]^\circ$ at
$T_i=8+0.25(i-1)~\mathrm{s}$. The initial joint angles differ from the tabulated
values by $[0.5,-0.4,0.4,-0.5,0.4,-0.4,0.3]^\circ$, giving
$\|e(0)\|=1.9\times10^{-2}~\mathrm{rad}$. The delay is $h=1~\mathrm{s}$, the
observation time is $T_o=1~\mathrm{s}$, hence $T=3~\mathrm{s}$, and
$R_h(\vartheta)=\sin^4(\pi(\vartheta-h)/h)$. The baseline gain is
$K_p=2I_7$, $K_d=2.8I_7$ (natural frequency $1.41~\mathrm{rad/s}$, damping
$0.99$) and $\mathrm{cond}(W_c)=7.4\times10^3$. The lumped disturbance is
applied at acceleration level,
$d_i=m_{e,i}(q_0)[0.02\sin(0.7t+0.3i)+0.01\cos(1.3t+0.2i)]$, so that
$\|\Delta_a\|\le3\times10^{-2}~\mathrm{rad/s^2}$ for every joint, which is the
form produced by inertial parameter errors. The observer uses $m_1=3$, $m_2=4$,
$\eta=0.3$, $\sigma=0.4$, and $\delta=0.05~\mathrm{s}$.

Six cases are compared: PD and PDF without disturbance, PD and PDF under the
disturbance, PDF with the observer, and the latter with a $3^\circ$ initial
error that forces the torques into the actuator limit.

\subsection{Results}
Table~\ref{tab:results} gives the metrics and Figs.~\ref{fig:main}--\ref{fig:obs}
the corresponding signals.

\emph{Nullification.} In the nominal case the delayed feedback annihilates the
error at the prescribed instant: the error over $[T,10]~\mathrm{s}$ is
$4.3\times10^{-5}~\mathrm{rad}$ against $1.4\times10^{-3}~\mathrm{rad}$ for PD,
and the measured settling time is $2.74~\mathrm{s}$ against $5.04~\mathrm{s}$.
The residual $1.7\times10^{-5}~\mathrm{rad}$ at $t=T$ is set by the Euler step
and the sampled delay buffer rather than by the design: halving $T_s$ halves
it. The periodic re-application of $K_c$ then keeps the error at the level of
one discretisation step, which is the practical counterpart of
Theorem~\ref{thm:null}.

\emph{Observer versus delayed feedback.} Without the observer the disturbed
equilibrium is $e\approx K_p^{-1}\Delta_a$, i.e., about
$1.5\times10^{-2}~\mathrm{rad}$ per joint, and the measured norms are
$2.4\times10^{-1}~\mathrm{rad}$ (PD) and $3.3\times10^{-1}~\mathrm{rad}$ (PDF).
Delayed feedback reshapes the transient but supplies no integral action, so it
cannot remove a persistent disturbance; the wrist joints dominate the norm
because their equivalent inertia is two to three decades below that of the
shoulder joints (Remark~\ref{rem:spread}). With the observer the estimation
error falls from $4.6\times10^{-1}$ to $1.5\times10^{-3}$ within
$T_o=1~\mathrm{s}$ (Fig.~\ref{fig:obs}) and the tracking error becomes
$9.3\times10^{-4}~\mathrm{rad}$ at $T$ and $1.5\times10^{-3}~\mathrm{rad}$
afterwards, with a peak torque of $12.4~\mathrm{N\,m}$.

\emph{Joint-level tracking.} Fig.~\ref{fig:joints} shows the joint positions
and rates together with their references. The executed signals follow the
reference during the manoeuvre and the residual mismatch is the error plotted
in Fig.~\ref{fig:errors}; the visible difference between PD and PDF appears
after $t=T$, where the delayed feedback suppresses the recursive lag of the
baseline while the torques remain smooth (Fig.~\ref{fig:torque}). No
high-frequency switching appears, which is the practical benefit of a
continuous delayed-feedback law over terminal sliding-mode designs.

\emph{Saturation.} A $3^\circ$ initial error demands about $37~\mathrm{N\,m}$,
so the limit is active for $3.6\%$ of the steps; the transient error reaches
$7.3\times10^{-2}~\mathrm{rad}$ and returns to the unsaturated value once the
error re-enters the actuator envelope. There is no windup because the law
\eqref{eq:law} has no integrator, but the error entering the design window is
larger, so the one-period map is only approximately annihilated during that
window, which is exactly the case covered by \eqref{eq:eps_bound}. An integral
or sliding-mode variant would instead require the fixed-time anti-saturation
compensator of~\cite{dou2025fractional}.

\emph{Base reaction.} The arm motion displaces the base by $6~\mathrm{cm}$ over
the manoeuvre (Fig.~\ref{fig:main}) with a monotone drift and no oscillatory
coupling, as predicted by momentum conservation. The displacement, together
with the associated attitude drift, is the sensor-frame bias that a task-space
implementation must absorb.

\begin{table}[!t]
\centering
\caption{Metrics. $\|e(T)\|$ is the error at $T=3~\mathrm{s}$ and
$\max_{t\ge T}\|e\|$ the residual afterwards; the saturation ratio is the
fraction of steps with $|\tau_i|=20~\mathrm{N\,m}$.}
\label{tab:results}
\scriptsize
\setlength{\tabcolsep}{1.5pt}
\begin{tabular}{lccccc}
\toprule
Case & $\|e(T)\|$ [rad] & $\max_{t\ge T}\|e\|$ & $\max|\tau_i|$ [N\,m] & sat. [\%] & $t_s$ [s]\\
\midrule
PD, nominal      & $1.60\times10^{-5}$ & $1.37\times10^{-3}$ & 10.84 & 0.0 & 5.04\\
PDF, nominal     & $1.65\times10^{-5}$ & $4.25\times10^{-5}$ & 12.19 & 0.0 & 2.74\\
PD, disturbed    & $2.39\times10^{-1}$ & $2.64\times10^{-1}$ & 11.94 & 0.0 & --\\
PDF, disturbed   & $3.31\times10^{-1}$ & $3.31\times10^{-1}$ & 20.00 & 7.5 & --\\
PDF+PTDO         & $9.28\times10^{-4}$ & $1.47\times10^{-3}$ & 12.38 & 0.0 & --\\
PDF+PTDO, $3^\circ$ & $9.28\times10^{-4}$ & $7.35\times10^{-2}$ & 20.00 & 3.6 & --\\
\bottomrule
\end{tabular}
\end{table}

\begin{figure*}[!t]
\centering
\subfloat[Base reaction and end-effector path\label{fig:motion}]{%
\includegraphics[width=0.30\textwidth]{figures/free_floating_trajectory.pdf}}
\hfill
\subfloat[Tracking-error norm\label{fig:errors}]{%
\includegraphics[width=0.30\textwidth]{figures/tracking_errors.pdf}}
\hfill
\subfloat[Control torque\label{fig:torque}]{%
\includegraphics[width=0.30\textwidth]{figures/control_torque.pdf}}
\caption{Free-floating motion and joint-space tracking under PDF+PTDO.
(a) momentum-induced base drift and end-effector path in the base frame;
(b) error norm for PD, PDF, and PDF+PTDO; (c) joint torques with the
$\pm20~\mathrm{N\,m}$ limits (dashed).}
\label{fig:main}
\end{figure*}

\begin{figure*}[!t]
\centering
\subfloat[Estimated disturbance\label{fig:est}]{%
\includegraphics[width=0.34\textwidth]{figures/observer_estimates.pdf}}
\hfill
\subfloat[Observation error\label{fig:oberr}]{%
\includegraphics[width=0.34\textwidth]{figures/observer_error.pdf}}
\caption{Prescribed-time observer: the estimate converges at $T_o=1~\mathrm{s}$
and the residual afterwards is of order $10^{-3}$, which dominates the
practical bound.}
\label{fig:obs}
\end{figure*}

\begin{figure*}[!t]
\centering
\subfloat[Joint positions\label{fig:jointpos}]{%
\includegraphics[width=0.42\textwidth]{figures/joint_position_tracking.pdf}}
\hfill
\subfloat[Joint rates\label{fig:jointvel}]{%
\includegraphics[width=0.42\textwidth]{figures/joint_velocity_tracking.pdf}}
\caption{Joint-space tracking under PDF+PTDO: executed (solid) and desired
(dashed) joint positions and rates.}
\label{fig:joints}
\end{figure*}

\section{Conclusion}\label{sec:conclusion}
A free-floating space manipulator is reduced exactly to a controllable error
channel, observed by a prescribed-time disturbance observer, and stabilised by
smooth periodic delayed feedback. The one-period map of the error channel is
annihilated in the ideal case, so the settling time is a design parameter; with
observation error, model mismatch, actuator limits, and sampling, an explicit
practical bound follows. On the three-dimensional model of a
seven-degree-of-freedom manipulator with published parameters and
$\pm20~\mathrm{N\,m}$ actuators, the delayed feedback reduces the post-settling
error by more than an order of magnitude and halves the settling time, while
the observer is what rejects a persistent disturbance. The study also exposes
the binding constraint: the equivalent inertia spans three decades, so a single
delay is ill matched to the fastest joint modes and the Gramian conditioning,
rather than the control law, sets the design margin.

Two modelling assumptions delimit the result: the momentum is zero, and the
plant is rigid apart from the actuator saturation. Nonzero momentum adds a
constant input channel that the observer would have to absorb, and appendage
flexibility introduces high-frequency dynamics that the delay buffer cannot
nullify; both would raise the residual in \eqref{eq:delta_r} and hence the
practical bound. The example also shows that the binding constraint is the
actuator envelope rather than the convergence mechanism: with
$\pm20~\mathrm{N\,m}$, a $1^\circ$ error on the shoulder joints already demands
about $11~\mathrm{N\,m}$.

Three extensions follow directly. A per-channel or multi-rate delay, with the
period map evaluated blockwise, would let the slow and fast joints be nullified
in their own time scales; the condition number of $W_c$ would then become a
design objective. A saturation-aware restart of the design window, or a
feasibility condition mapping the initial error, the torque limit, and the
delay into an admissible set, would restore the exact guarantee in the
torque-limited regime. Finally, closing the loop through the generalized
Jacobian with a reaction null space, and validating the delay buffer on
hardware, would connect the present joint-space result to a task-space
on-orbit servicing controller.


\appendices
\section{Proofs}\label{app:proofs}

\subsection{Positive definiteness of the Gramian}
For any $x\in\mathbb R^{2n}$,
\begin{equation}
 x^\top W_cx=\int_h^{2h}R_h(s)\,\|B^\top e^{-A_c^\top s}x\|^2ds\ge0 .
\end{equation}
If the integral vanishes, continuity gives $B^\top e^{-A_c^\top s}x=0$ on an
interval contained in the set where $R_h>0$. The function
$s\mapsto B^\top e^{-A_c^\top s}x$ is analytic, hence it vanishes for all $s$,
and differentiating at any point gives $B^\top(A_c^\top)^kx=0$ for every
$k\ge0$. Controllability of $(A_c,B)$ then implies $x=0$, so $W_c\succ0$.

\subsection{Observer error system}
Subtracting \eqref{eq:ptdo} from \eqref{eq:channel} with
$\varepsilon_1=\chi-z_1$ and $\varepsilon_2=\Delta_a-z_2$ yields
\begin{equation}
 \dot\varepsilon_1=\varepsilon_2-\ell_1(t)\phi_1(\varepsilon_1/\sigma),
 \qquad
 \dot\varepsilon_2=-\ell_2(t)\phi_2(\varepsilon_1/\sigma)+\dot\Delta_a .
 \label{eq:obs_err}
\end{equation}
Let $s=-\ln(1-t/T_o)$, so $ds/dt=1/(T_o-t)$. Since
$\ell_1=m_1\,ds/dt$ and $\ell_2=m_2\,(ds/dt)^2$, the unperturbed observer is a
homogeneous double integrator with constant gains in the stretched time.
Because $\phi_1$ and $\phi_2$ are odd, continuous, and
strictly increasing, and since the powers satisfy
$1\mp\eta=2(1\mp\eta/2)-1$, the unperturbed field of \eqref{eq:obs_err} is
homogeneous with weights $(1,1\mp\eta/2)$ and has negative degree on both
branches; hence it is finite-time stable in $s$, which corresponds to
convergence at $t=T_o$ in the original time. With
$\|\dot\Delta_a\|\le L_d$ the perturbation is bounded, so the perturbed
solution enters every neighbourhood of the origin before $T_o$, and strong
invariance of $\{\varepsilon_1=\varepsilon_2=0\}$ gives
$z_2(t)=\Delta_a(t)$ for $t\ge T_o$. The regularised implementation freezes the
gains at $T_o-\delta$, which replaces exactness by the residual
$O(L_d\delta^2/m_2)$ used in \eqref{eq:delta_r}.

\subsection{Explicit residual}
Combining \eqref{eq:reduced} with the implemented torque \eqref{eq:law} and
writing $\tilde M=M_e^{-1}\hat M_e-I_n$ for the relative model error gives
\begin{equation}
 \ddot e=\nu+\tilde M\bigl(\ddot q_d+\nu\bigr)
 +M_e^{-1}\bigl[(\hat C_e-C_e)\dot q+(d-\hat d)\bigr],
\end{equation}
so that the residual acceleration in \eqref{eq:error_channel} is bounded by
\begin{equation}
 \|\Delta_r\|\le\|\tilde M\|\,(\|\ddot q_d+\nu\|)
 +\underline m^{-1}\bigl(\bar d_o+\bar\Delta_m+\bar\Delta_\tau\bigr),
\end{equation}
with $\bar d_o$ the observation residual, $\bar\Delta_m$ the torque-level model
error, and $\bar\Delta_\tau$ the torque clipping. This is the estimate used in
\eqref{eq:bound}--\eqref{eq:eps_bound}.


\bibliographystyle{IEEEtran}
\bibliography{refs/local}

% IEEE TAES publishes short author biographies. Uncomment and complete for
% the final files; they are omitted from the review version.
% % Author biographies (IEEE TAES style). Fill in and uncomment the
% % Author biographies (IEEE TAES style). Fill in and uncomment the
% % Author biographies (IEEE TAES style). Fill in and uncomment the
% \input{taes/bios} line in paper_taes.tex for the final files.

\begin{IEEEbiographynophoto}{Author Name}
received the B.S. degree in XXX from XXX University, XXX, China, in 20XX, and
the Ph.D. degree in XXX from XXX University, XXX, China, in 20XX. He/She is
currently a XXX with the Department of XXX, XXX University. His/Her research
interests include space robotics, nonlinear control, and time-delay systems.
\end{IEEEbiographynophoto}

\begin{IEEEbiographynophoto}{Author Name}
received the Ph.D. degree in XXX from XXX University, XXX, China, in 20XX. He
is currently a Professor with the Department of XXX, XXX University. His
research interests include spacecraft dynamics and control and on-orbit
servicing.
\end{IEEEbiographynophoto}
 line in paper_taes.tex for the final files.

\begin{IEEEbiographynophoto}{Author Name}
received the B.S. degree in XXX from XXX University, XXX, China, in 20XX, and
the Ph.D. degree in XXX from XXX University, XXX, China, in 20XX. He/She is
currently a XXX with the Department of XXX, XXX University. His/Her research
interests include space robotics, nonlinear control, and time-delay systems.
\end{IEEEbiographynophoto}

\begin{IEEEbiographynophoto}{Author Name}
received the Ph.D. degree in XXX from XXX University, XXX, China, in 20XX. He
is currently a Professor with the Department of XXX, XXX University. His
research interests include spacecraft dynamics and control and on-orbit
servicing.
\end{IEEEbiographynophoto}
 line in paper_taes.tex for the final files.

\begin{IEEEbiographynophoto}{Author Name}
received the B.S. degree in XXX from XXX University, XXX, China, in 20XX, and
the Ph.D. degree in XXX from XXX University, XXX, China, in 20XX. He/She is
currently a XXX with the Department of XXX, XXX University. His/Her research
interests include space robotics, nonlinear control, and time-delay systems.
\end{IEEEbiographynophoto}

\begin{IEEEbiographynophoto}{Author Name}
received the Ph.D. degree in XXX from XXX University, XXX, China, in 20XX. He
is currently a Professor with the Department of XXX, XXX University. His
research interests include spacecraft dynamics and control and on-orbit
servicing.
\end{IEEEbiographynophoto}


\end{document}
